\documentclass[12pt,a4paper]{article}

% ── Packages ──────────────────────────────────────────────────────────────────
\usepackage[margin=2.5cm]{geometry}
\usepackage{graphicx}
\usepackage{booktabs}
\usepackage{array}
\usepackage{longtable}
\usepackage{hyperref}
\usepackage{xcolor}
\usepackage{amsmath}
\usepackage{enumitem}
\usepackage{titlesec}
\usepackage{fancyhdr}
\usepackage{listings}
\usepackage{caption}
\usepackage{subcaption}
\usepackage{tabularx}
\usepackage{multirow}
\usepackage{float}
\usepackage{parskip}
\usepackage{mdframed}
\usepackage{fontenc}
\usepackage{inputenc}
\usepackage{microtype}
\usepackage{url}
\usepackage{tcolorbox}
\tcbuselibrary{skins,breakable}

% ── Colours ───────────────────────────────────────────────────────────────────
\definecolor{primary}{RGB}{30, 90, 160}
\definecolor{accent}{RGB}{220, 50, 50}
\definecolor{lightgray}{RGB}{245, 245, 245}
\definecolor{codebg}{RGB}{248, 248, 248}
\definecolor{churnred}{RGB}{200, 40, 40}
\definecolor{retaingreen}{RGB}{30, 140, 70}

% ── Hyperlinks ────────────────────────────────────────────────────────────────
\hypersetup{
    colorlinks=true,
    linkcolor=primary,
    urlcolor=primary,
    citecolor=primary,
    pdftitle={Customer Churn Prediction -- Technical Report},
    pdfauthor={Pushpendra S. Parihar, Rohan Singh, Saksham Sharma}
}

% ── Section formatting ────────────────────────────────────────────────────────
\titleformat{\section}{\large\bfseries\color{primary}}{{\thesection}}{1em}{}[\titlerule]
\titleformat{\subsection}{\normalsize\bfseries\color{primary}}{\thesubsection}{1em}{}
\titleformat{\subsubsection}{\normalsize\itshape}{\thesubsubsection}{1em}{}

% ── Header / Footer ───────────────────────────────────────────────────────────
\pagestyle{fancy}
\fancyhf{}
\rhead{\small\color{gray} Customer Churn Prediction}
\lhead{\small\color{gray} Technical Report --- Milestone 1}
\rfoot{\small\thepage}
\lfoot{\small\color{gray} Gen AI Assignment}
\renewcommand{\headrulewidth}{0.4pt}

% ── Code listing style ────────────────────────────────────────────────────────
\lstset{
    backgroundcolor=\color{codebg},
    basicstyle=\ttfamily\small,
    breaklines=true,
    frame=single,
    rulecolor=\color{gray!40},
    xleftmargin=10pt,
    xrightmargin=10pt,
    captionpos=b
}

% ── Metric box ────────────────────────────────────────────────────────────────
\newtcolorbox{metricbox}[1]{
    colback=lightgray,
    colframe=primary,
    title=#1,
    fonttitle=\bfseries\small,
    left=6pt, right=6pt, top=4pt, bottom=4pt
}

% ═══════════════════════════════════════════════════════════════════════════════
\begin{document}
% ═══════════════════════════════════════════════════════════════════════════════

% ── Title Page ────────────────────────────────────────────────────────────────
\begin{titlepage}
    \centering
    \vspace*{1cm}
    {\color{primary}\rule{\linewidth}{2pt}}
    \vspace{0.5cm}

    {\LARGE\bfseries\color{primary} Customer Churn Prediction \\ \& Retention Strategy}

    \vspace{0.3cm}
    {\large\color{gray} Technical Report --- Milestone 1}

    \vspace{0.3cm}
    {\color{primary}\rule{\linewidth}{0.5pt}}

    \vspace{1.5cm}
    \begin{tcolorbox}[colback=lightgray, colframe=primary!40, width=0.85\linewidth]
        \centering
        \textbf{Live Deployment:}\\[4pt]
        \href{https://customer-churn-predictiongit-mid-sem-milestone-1.streamlit.app/}
             {\ttfamily\small customer-churn-predictiongit-mid-sem-milestone-1.streamlit.app}
    \end{tcolorbox}

    \vspace{1cm}
    \begin{tabular}{rl}
        \textbf{Authors}     & Pushpendra S. Parihar, Rohan Singh \& Saksham Sharma \\[4pt]
        \textbf{GitHub}      & \href{https://github.com/makeprodigy/customer-churn-prediction}{makeprodigy/customer-churn-prediction} \\[4pt]
        \textbf{Course}      & Generative AI --- Mid-Semester Project (Project 5) \\[4pt]
        \textbf{Submitted}   & 2 March 2026 \\[4pt]
        \textbf{Tech Stack}  & Python 3 · Scikit-Learn · Streamlit · Pandas · Seaborn \\
    \end{tabular}

    \vfill
    {\color{primary}\rule{\linewidth}{2pt}}
\end{titlepage}

% ── Table of Contents ─────────────────────────────────────────────────────────
\tableofcontents
\newpage

% ═══════════════════════════════════════════════════════════════════════════════
\section{Problem Statement}
% ═══════════════════════════════════════════════════════════════════════════════

Customer churn --- the voluntary discontinuation of service by a subscriber --- is one of the most economically damaging phenomena in the telecommunications industry. Acquiring a new customer is estimated to cost five to seven times more than retaining an existing one, and industry research suggests that even a \textbf{5\% reduction in churn can increase profitability by 25--95\%}.

\subsection{Challenge Definition}

Given a rich set of 19 customer attributes spanning demographics, service subscriptions, account details, and billing behaviour, the challenge is to build a reliable binary classification system that answers:

\begin{quote}
\itshape
\textbf{``Will this customer churn within the next billing cycle?''}
\end{quote}

More specifically, the system must:
\begin{enumerate}[leftmargin=2em]
    \item \textbf{Predict} whether a given customer will churn (\texttt{Churn = Yes}) or remain (\texttt{Churn = No}).
    \item \textbf{Quantify} the probability of churn to enable risk-tiered retention prioritisation.
    \item \textbf{Operate} in two modes: real-time single-customer inference and bulk batch processing.
    \item \textbf{Generalise} robustly to unseen customer profiles without overfitting to training noise.
\end{enumerate}

\subsection{Business Impact}

Accurate churn prediction allows the business team to:
\begin{itemize}[leftmargin=2em]
    \item Target high-risk customers with personalised retention offers \emph{before} they leave.
    \item Prioritise customer success resources efficiently.
    \item Identify systemic issues driving disengagement (e.g., contract type, internet service quality).
\end{itemize}

% ═══════════════════════════════════════════════════════════════════════════════
\section{Data Description}
% ═══════════════════════════════════════════════════════════════════════════════

\subsection{Dataset Source}

The project uses the \textbf{IBM Telco Customer Churn} dataset, a widely-used public benchmark for binary churn classification. The dataset represents customers of a fictional telecom provider and is loaded directly from \texttt{customer\_churn\_datasest.csv}.

\subsection{Dataset Statistics}

\begin{table}[H]
\centering
\caption{High-Level Dataset Summary}
\begin{tabular}{ll}
    \toprule
    \textbf{Property} & \textbf{Value} \\
    \midrule
    Total records     & 7,043 customers \\
    Total features    & 20 (1 ID + 19 predictors) \\
    Target variable   & \texttt{Churn} (binary: Yes / No) \\
    Churn rate        & 26.54\% (1,869 churned / 5,174 retained) \\
    Missing values    & 11 rows in \texttt{TotalCharges} (blank strings → NaN) \\
    \bottomrule
\end{tabular}
\end{table}

\subsection{Feature Inventory}

\begin{table}[H]
\centering
\caption{Feature Descriptions}
\begin{tabularx}{\textwidth}{llX}
    \toprule
    \textbf{Feature} & \textbf{Type} & \textbf{Description} \\
    \midrule
    \texttt{customerID}        & Categorical & Unique identifier (dropped before modelling) \\
    \texttt{gender}            & Categorical & Male / Female \\
    \texttt{SeniorCitizen}     & Binary      & 1 if customer is 65+, else 0 \\
    \texttt{Partner}           & Binary      & Has a partner (Yes/No) \\
    \texttt{Dependents}        & Binary      & Has dependents (Yes/No) \\
    \texttt{tenure}            & Numerical   & Months with the company (0--72) \\
    \texttt{PhoneService}      & Binary      & Has phone service \\
    \texttt{MultipleLines}     & Categorical & No / Yes / No phone service \\
    \texttt{InternetService}   & Categorical & DSL / Fiber optic / No \\
    \texttt{OnlineSecurity}    & Categorical & No / Yes / No internet service \\
    \texttt{OnlineBackup}      & Categorical & No / Yes / No internet service \\
    \texttt{DeviceProtection}  & Categorical & No / Yes / No internet service \\
    \texttt{TechSupport}       & Categorical & No / Yes / No internet service \\
    \texttt{StreamingTV}       & Categorical & No / Yes / No internet service \\
    \texttt{StreamingMovies}   & Categorical & No / Yes / No internet service \\
    \texttt{Contract}          & Categorical & Month-to-month / One year / Two year \\
    \texttt{PaperlessBilling}  & Binary      & Enrolled in paperless billing \\
    \texttt{PaymentMethod}     & Categorical & Electronic check / Mailed check / etc. \\
    \texttt{MonthlyCharges}    & Numerical   & Monthly bill amount (\$18.25--\$118.75) \\
    \texttt{TotalCharges}      & Numerical   & Cumulative spend (NaN-imputed) \\
    \texttt{Churn}             & \textcolor{churnred}{\textbf{Target}} & Yes (1) or No (0) \\
    \bottomrule
\end{tabularx}
\end{table}

\subsection{Class Imbalance Note}

The dataset exhibits moderate class imbalance with a 73.5\% / 26.5\% split. This was addressed during model evaluation by using \textbf{F1-Score} as the primary optimisation metric (rather than accuracy alone), which prevents the classifier from being biased toward predicting the majority class.

% ═══════════════════════════════════════════════════════════════════════════════
\section{Exploratory Data Analysis (EDA)}
% ═══════════════════════════════════════════════════════════════════════════════

\subsection{Numerical Feature Statistics}

\begin{table}[H]
\centering
\caption{Descriptive Statistics for Numerical Features}
\begin{tabular}{lrrrrrr}
    \toprule
    \textbf{Feature} & \textbf{Mean} & \textbf{Std} & \textbf{Min} & \textbf{25\%} & \textbf{Median} & \textbf{Max} \\
    \midrule
    \texttt{tenure}         & 32.37  & 24.56 & 0     & 9      & 29      & 72       \\
    \texttt{MonthlyCharges} & \$64.76 & \$30.09 & \$18.25 & \$35.50 & \$70.35 & \$118.75 \\
    \texttt{TotalCharges}   & \$2,283 & \$2,267 & \$18.80 & \$401  & \$1,397 & \$8,685  \\
    \bottomrule
\end{tabular}
\end{table}

\subsection{Churn by Key Categorical Features}

\begin{table}[H]
\centering
\caption{Churn Rate by Key Categorical Variables}
\begin{tabular}{llrr}
    \toprule
    \textbf{Feature} & \textbf{Value} & \textbf{Churn \%} & \textbf{Insight} \\
    \midrule
    \multirow{3}{*}{Contract}
        & Month-to-month  & \textcolor{churnred}{\textbf{42.7\%}} & Highest single driver \\
        & One year        & 11.3\% & Moderate retention \\
        & Two year        & \textcolor{retaingreen}{\textbf{2.8\%}}  & Strongest retention \\
    \midrule
    \multirow{3}{*}{InternetService}
        & Fiber optic     & \textcolor{churnred}{\textbf{41.9\%}} & High churn service tier \\
        & DSL             & 19.0\% & Lower churn \\
        & No internet     & \textcolor{retaingreen}{\textbf{7.4\%}}  & Very loyal \\
    \midrule
    \multirow{4}{*}{PaymentMethod}
        & Electronic check            & \textcolor{churnred}{\textbf{45.3\%}} & Disengaged payment \\
        & Mailed check                & 19.1\% & \\
        & Bank transfer (automatic)   & 16.7\% & \\
        & Credit card (automatic)     & \textcolor{retaingreen}{\textbf{15.2\%}} & Most loyal group \\
    \midrule
    \multirow{2}{*}{SeniorCitizen}
        & 1 (Senior)     & \textcolor{churnred}{\textbf{41.7\%}} & Vulnerable segment \\
        & 0 (Non-Senior) & 23.6\% & \\
    \midrule
    \multirow{2}{*}{PaperlessBilling}
        & Yes            & 33.6\% & Higher churn \\
        & No             & \textcolor{retaingreen}{\textbf{16.3\%}} & \\
    \midrule
    \multirow{2}{*}{Dependents}
        & No             & 31.3\% & No dependents = more mobile \\
        & Yes            & \textcolor{retaingreen}{\textbf{15.5\%}} & Families are stickier \\
    \bottomrule
\end{tabular}
\end{table}

\subsection{Numerical Features: Churn vs.\ Non-Churn Comparison}

\begin{table}[H]
\centering
\caption{Mean Values by Churn Status}
\begin{tabular}{lrrl}
    \toprule
    \textbf{Feature} & \textbf{Churned} & \textbf{Retained} & \textbf{Interpretation} \\
    \midrule
    \texttt{tenure}           & 17.98 mo  & 37.57 mo  & Churned customers leave ~2× earlier \\
    \texttt{MonthlyCharges}   & \$74.44   & \$61.27   & Churners pay more per month \\
    \texttt{TotalCharges}     & \$1,531   & \$2,555   & Less lifetime value (short tenure) \\
    \bottomrule
\end{tabular}
\end{table}

\subsection{Correlation with Churn}

Pearson correlations between numerical features and the binary churn label:

\begin{table}[H]
\centering
\caption{Pearson Correlation with \texttt{Churn}}
\begin{tabular}{lrl}
    \toprule
    \textbf{Feature} & \textbf{Correlation} & \textbf{Direction} \\
    \midrule
    \texttt{tenure}           & $-0.352$ & Longer tenure $\Rightarrow$ less churn \\
    \texttt{TotalCharges}     & $-0.199$ & Higher total spend $\Rightarrow$ less churn (proxy for tenure) \\
    \texttt{MonthlyCharges}   & $+0.193$ & Higher monthly bill $\Rightarrow$ more churn \\
    \texttt{SeniorCitizen}    & $+0.151$ & Being a senior $\Rightarrow$ more churn \\
    \bottomrule
\end{tabular}
\end{table}

\subsection{Key EDA Takeaways}

\begin{tcolorbox}[colback=lightgray, colframe=primary, title={\textbf{EDA Summary --- Top Churn Drivers}}]
\begin{enumerate}[leftmargin=1.5em, itemsep=2pt]
    \item \textbf{Contract type} is the dominant predictor: month-to-month customers churn at 15$\times$ the rate of two-year contract holders.
    \item \textbf{Fiber optic internet} paradoxically shows high churn (41.9\%), suggesting a cost/expectation mismatch.
    \item \textbf{Electronic check} users are the most likely to churn (45.3\%) --- a disengaged payment method signals disengaged customers.
    \item \textbf{Tenure} is the strongest negative predictor ($r = -0.352$): the first 12 months are critical for retention.
    \item \textbf{Senior citizens} without partners/dependents represent the most at-risk demographic segment.
\end{enumerate}
\end{tcolorbox}

% ═══════════════════════════════════════════════════════════════════════════════
\section{Methodology}
% ═══════════════════════════════════════════════════════════════════════════════

\subsection{System Architecture}

The system is structured as a two-component pipeline:

\begin{enumerate}[leftmargin=2em]
    \item \textbf{Offline Training} (\texttt{src/preprocess.py} + \texttt{src/train.py}): Builds and serialises scikit-learn \texttt{Pipeline} objects to \texttt{.joblib} files.
    \item \textbf{Online Inference} (\texttt{app.py}): Loads \texttt{.joblib} files and serves predictions via a Streamlit web application with two modes: single-customer and batch CSV.
\end{enumerate}

\subsection{Preprocessing Pipeline}

To ensure a \textbf{leak-proof} preprocessing strategy (i.e., fitting scalers only on training data), all transformations are encapsulated inside a \texttt{ColumnTransformer} that becomes part of each model's \texttt{Pipeline}:

\begin{lstlisting}[language=Python, caption={Preprocessing Pipeline (src/preprocess.py)}]
numeric_transformer = Pipeline([
    ('imputer', SimpleImputer(strategy='median')),
    ('scaler',  StandardScaler())
])

categorical_transformer = Pipeline([
    ('imputer', SimpleImputer(strategy='most_frequent')),
    ('onehot',  OneHotEncoder(handle_unknown='ignore'))
])

preprocessor = ColumnTransformer([
    ('num', numeric_transformer, numeric_features),    # tenure, MonthlyCharges, TotalCharges
    ('cat', categorical_transformer, categorical_features)   # 16 categorical columns
])
\end{lstlisting}

\subsubsection{Design Decisions}
\begin{itemize}[leftmargin=2em]
    \item \textbf{Median imputation} for \texttt{TotalCharges}: robust to the skewed distribution caused by early-churning customers.
    \item \textbf{StandardScaler} for numerical features: makes regularised models (LR, MLP) converge faster and improves coefficient interpretability.
    \item \textbf{One-Hot Encoding with \texttt{handle\_unknown='ignore'}}: prevents errors when serving predictions on new customers with unseen category values.
    \item \textbf{Train/test split}: 80\% / 20\% stratified by random seed 42 (5,634 train / 1,409 test).
\end{itemize}

\subsection{Classifier Selection \& Rationale}

\begin{table}[H]
\centering
\caption{Model Selection Rationale}
\begin{tabularx}{\textwidth}{llX}
    \toprule
    \textbf{Model} & \textbf{Library Class} & \textbf{Rationale} \\
    \midrule
    Logistic Regression
        & \texttt{LogisticRegression}
        & Gold-standard interpretable baseline for binary classification. Linear decision boundary is suitable for well-engineered tabular features. Coefficients directly indicate feature influence.\\[6pt]
    Decision Tree
        & \texttt{DecisionTreeClassifier}
        & Captures non-linear feature interactions and provides human-readable classification rules. Useful for explaining individual predictions to non-technical stakeholders. \\[6pt]
    MLP (Neural Net)
        & \texttt{MLPClassifier}
        & Multi-layer perceptron as a deep-learning baseline. Capable of learning complex feature interactions automatically. Serves as an upper bound on achievable accuracy without manual feature engineering. \\
    \bottomrule
\end{tabularx}
\end{table}

Each classifier is wrapped in a full \texttt{sklearn.pipeline.Pipeline}:
\[
\texttt{Pipeline} = [\; \underbrace{\texttt{ColumnTransformer}}_{\text{Preprocessing}} \;\rightarrow\; \underbrace{\texttt{Classifier}}_{\text{Model}} \;]
\]
This ensures preprocessing and inference are always applied consistently, whether during training, cross-validation, or production scoring.

% ═══════════════════════════════════════════════════════════════════════════════
\section{Evaluation}
% ═══════════════════════════════════════════════════════════════════════════════

\subsection{Metrics Used}

Given the class imbalance (26.5\% positive class), multiple metrics were used to obtain a complete picture of model performance:

\begin{itemize}[leftmargin=2em]
    \item \textbf{Accuracy}: Overall fraction of correct predictions.
    \item \textbf{Precision}: Of all customers predicted to churn, how many actually did? (Minimises false alarms.)
    \item \textbf{Recall}: Of all customers who actually churned, how many were caught? (Minimises missed churners.)
    \item \textbf{F1 Score}: Harmonic mean of precision and recall --- the primary optimisation metric, selected because it balances both objectives under class imbalance.
\end{itemize}

\subsection{Holdout Test Results}

All models were evaluated on a held-out test set of \textbf{1,409 customers} (20\% of the full dataset):

\begin{table}[H]
\centering
\caption{Model Evaluation Results on 20\% Holdout Test Set}
\begin{tabular}{lrrrr}
    \toprule
    \textbf{Model} & \textbf{Accuracy} & \textbf{Precision} & \textbf{Recall} & \textbf{F1 Score} \\
    \midrule
    \textbf{Logistic Regression} & \textbf{81.97\%} & \textbf{68.42\%} & 59.25\% & \textbf{63.51\%} \\
    Decision Tree                & 79.99\%          & 61.88\%          & \textbf{63.54\%} & 62.70\% \\
    MLP (Neural Net)             & 78.57\%          & 61.41\%          & 51.21\% & 55.85\% \\
    \bottomrule
\end{tabular}
\end{table}

\subsection{Result Interpretation}

\begin{tcolorbox}[colback=lightgray, colframe=retaingreen, title={\textbf{Winner: Logistic Regression}}]
Logistic Regression achieves the highest accuracy (81.97\%) and F1-Score (63.51\%), making it the recommended model for production deployment. Its linear coefficients confirm what EDA revealed: \emph{month-to-month contracts}, \emph{fiber optic internet}, and \emph{high monthly charges} are the strongest positive drivers of churn.
\end{tcolorbox}

\begin{itemize}[leftmargin=2em]
    \item \textbf{Decision Tree} achieves the highest Recall (63.54\%), meaning it catches more actual churners --- useful if the business priority is to \emph{minimise missed churners} at the cost of some false alarms.
    \item \textbf{MLP} underperforms on Recall (51.21\%), suggesting the dataset size (7K rows) is likely insufficient to fully leverage a neural network's representational capacity without more aggressive regularisation or data augmentation.
\end{itemize}

% ═══════════════════════════════════════════════════════════════════════════════
\section{Optimisation}
% ═══════════════════════════════════════════════════════════════════════════════

\subsection{Hyperparameter Tuning Strategy}

All three models were optimised using \textbf{GridSearchCV with 3-fold cross-validation}, scoring on F1 to handle class imbalance correctly:

\begin{lstlisting}[language=Python, caption={GridSearchCV Setup (src/train.py)}]
grid_search = GridSearchCV(
    estimator=full_pipeline,
    param_grid=config['params'],
    scoring='f1',      # optimise for F1, not accuracy
    cv=3,
    n_jobs=-1          # parallelise across all CPU cores
)
grid_search.fit(X_train, y_train)
best_pipeline = grid_search.best_estimator_
\end{lstlisting}

\subsection{Search Spaces \& Best Parameters}

\begin{table}[H]
\centering
\caption{GridSearchCV Parameter Grids and Optimal Results}
\begin{tabularx}{\textwidth}{llXl}
    \toprule
    \textbf{Model} & \textbf{Parameter} & \textbf{Candidates Searched} & \textbf{Best Value} \\
    \midrule
    \multirow{2}{*}{Logistic Regression}
        & \texttt{C} (inverse regularisation)
        & \{0.1, 1.0, 10.0\}
        & \textbf{10.0} \\
        & \texttt{solver}
        & \{lbfgs, liblinear\}
        & \textbf{liblinear} \\
    \midrule
    \multirow{2}{*}{Decision Tree}
        & \texttt{max\_depth}
        & \{3, 5, 10, None\}
        & \textbf{5} \\
        & \texttt{min\_samples\_leaf}
        & \{1, 5, 10\}
        & \textbf{10} \\
    \midrule
    \multirow{2}{*}{MLP}
        & \texttt{hidden\_layer\_sizes}
        & \{(50,), (100,)\}
        & \textbf{(50,)} \\
        & \texttt{alpha} (L2 penalty)
        & \{0.0001, 0.001\}
        & \textbf{0.0001} \\
    \bottomrule
\end{tabularx}
\end{table}

\subsection{Optimisation Impact \& Design Rationale}

\begin{itemize}[leftmargin=2em]
    \item \textbf{Logistic Regression}: The best \texttt{C=10} corresponds to mild regularisation, implying the dataset is clean and most features are informative. \texttt{liblinear} was selected over \texttt{lbfgs} — it performs better on smaller, sparse one-hot-encoded feature spaces.

    \item \textbf{Decision Tree}: Constraining \texttt{max\_depth=5} and \texttt{min\_samples\_leaf=10} \emph{significantly} mitigated overfitting. Without these constraints, an unconstrained tree achieved 100\% training accuracy but scored below 75\% on the test set.

    \item \textbf{MLP}: The smaller \texttt{(50,)} hidden layer (vs. \texttt{(100,)}) performed better, confirming the dataset is not complex enough to warrant a larger network. A lower \texttt{alpha} (less regularisation) suggests the MLP benefits from more model freedom given the limited parameters.

    \item \textbf{Cross-validation}: 3-fold CV across all 18/36/4 parameter combinations (totalling 18 + 36 + 12 = 66 fits) ensures robust generalisation estimates unaffected by the particular random state of a single train/validation split.
\end{itemize}

% ═══════════════════════════════════════════════════════════════════════════════
\section{Team Contribution}
% ═══════════════════════════════════════════════════════════════════════════════

\subsection{Contributor Overview}

This project was developed collaboratively by three contributors. Git commit responsibilities are split in a \textbf{4:4:2 ratio} (Pushpendra~40\%, Rohan~40\%, Saksham~20\%), with additional non-code deliverables assigned as described below.

\begin{table}[H]
\centering
\caption{Contributor Summary}
\begin{tabular}{llrp{5.5cm}}
    \toprule
    \textbf{Contributor} & \textbf{GitHub Handle} & \textbf{Git Share} & \textbf{Key Responsibilities} \\
    \midrule
    Pushpendra S. Parihar & \texttt{makeprodigy}
        & 40\%
        & ML pipeline, Streamlit app, hyperparameter optimisation, cloud deployment, \textbf{technical report (co-author)} \\
    Rohan Singh & \texttt{r0hansng}
        & 40\%
        & Data sourcing, EDA, model evaluation analysis, \textbf{project demo video} \\
    Saksham Sharma & \texttt{Saksham-9898}
        & 20\%
        & Model evaluation review, testing, documentation feedback, \textbf{technical report (co-author)} \\
    \bottomrule
\end{tabular}
\end{table}

\subsection{Individual Responsibility Breakdown}

\begin{table}[H]
\centering
\caption{Role-Based Contribution Split}
\begin{tabularx}{\textwidth}{lXXX}
    \toprule
    \textbf{Task Area} & \textbf{Pushpendra S. Parihar} & \textbf{Rohan Singh} & \textbf{Saksham Sharma} \\
    \midrule
    Project Setup \& Scaffold
        & Initialised repo, gitignore, requirements
        & Reviewed setup
        & Reviewed setup \\
    Data Pipeline
        & Authored \texttt{preprocess.py}: imputation, scaling, one-hot encoding
        & Sourced dataset, validated columns \& missing values
        & Cross-checked encoding logic \\
    Model Training
        & Authored \texttt{train.py}: LR, DT, MLP with GridSearchCV
        & Reviewed model configs, analysed metrics
        & Validated evaluation loop outputs \\
    Streamlit Application
        & Built full dashboard (\texttt{app.py}): single \& batch modes, charts, session state
        & User testing \& feedback across UI iterations
        & Functional testing \& edge-case validation \\
    Cloud Deployment
        & Deployed to Streamlit Cloud; committed \texttt{.joblib} artefacts
        & Verified live URL \& behaviour
        & — \\
    Hyperparameter Opt.
        & Implemented GridSearchCV; regenerated models
        & Compared before/after metrics
        & Documented tuning results \\
    Documentation (README)
        & Wrote \& iterated README (3 commits)
        & Reviewed content
        & Provided feedback \& proofreading \\
    \textbf{Project Demo Video}
        & —
        & \textbf{Produced, recorded \& edited full demo video}
        & — \\
    \textbf{Technical Report}
        & \textbf{LaTeX authoring \& structure}
        & —
        & \textbf{Content review, EDA write-up \& proofreading} \\
    \bottomrule
\end{tabularx}
\end{table}


\subsection{Effort Distribution Summary}

The demo video is available at: \href{https://drive.google.com/file/d/1u08SJbGc92HiiyLoe42iHyELQlHl03aQ/view?usp=sharing}{\texttt{Google Drive (Demo Video)}}.

\begin{table}[H]
\centering
\caption{Overall Contribution Split (4:4:2 Ratio)}
\begin{tabular}{llrp{6cm}}
    \toprule
    \textbf{Contributor} & \textbf{Git Share} & \textbf{\%} & \textbf{Extra Deliverables} \\
    \midrule
    Pushpendra S. Parihar & Phases 1--7 (lead) & 40\% & Technical report (LaTeX) --- co-author \\
    Rohan Singh           & Phases 2--7 (EDA, testing) & 40\% & Project demo video (\href{https://drive.google.com/file/d/1u08SJbGc92HiiyLoe42iHyELQlHl03aQ/view?usp=sharing}{link}) \\
    Saksham Sharma        & Phases 1, 6--7 (review) & 20\% & Technical report (LaTeX) --- co-author \\
    \midrule
    \textbf{Total} & & \textbf{100\%} & \\
    \bottomrule
\end{tabular}
\end{table}

% ═══════════════════════════════════════════════════════════════════════════════
\section*{Conclusion}
\addcontentsline{toc}{section}{Conclusion}
% ═══════════════════════════════════════════════════════════════════════════════

This report documents the end-to-end development of a customer churn prediction system for Milestone 1 of Project 5. Key outcomes:

\begin{itemize}[leftmargin=2em]
    \item A \textbf{leak-proof scikit-learn pipeline} was engineered combining median/mode imputation, standardisation, and one-hot encoding, all encapsulated within each model's \texttt{Pipeline} object.
    \item \textbf{Three classifiers} (Logistic Regression, Decision Tree, MLP) were trained and tuned via GridSearchCV with 3-fold cross-validation.
    \item \textbf{Logistic Regression} emerged as the best model with \textbf{81.97\% accuracy} and \textbf{63.51\% F1-Score}.
    \item A \textbf{fully deployed Streamlit application} provides real-time single-customer and bulk batch churn predictions, accessible at the public URL listed on the title page.
    \item The project was delivered as a \textbf{collaborative team effort} across 15 well-structured commits: Pushpendra S. Parihar (40\%, code, deployment, report), Rohan Singh (40\%, EDA, testing, demo video), and Saksham Sharma (20\%, evaluation review, documentation, report).
\end{itemize}

\subsection*{Milestone 2 Outlook}
The next milestone will extend this system into an \textbf{agentic AI retention strategist} using LangGraph and RAG, capable of reasoning about detected churn risk and autonomously proposing personalised retention strategies.

% ═══════════════════════════════════════════════════════════════════════════════
\newpage
\section*{References}
\addcontentsline{toc}{section}{References}
% ═══════════════════════════════════════════════════════════════════════════════

\begin{enumerate}[leftmargin=2em, itemsep=6pt]
    \item IBM. (2019). \textit{Telco Customer Churn Dataset}. IBM Sample Data Sets. \\
          \url{https://www.kaggle.com/datasets/blastchar/telco-customer-churn}

    \item Pedregosa, F., Varoquaux, G., Gramfort, A., et al. (2011). Scikit-learn: Machine Learning in Python.
          \textit{Journal of Machine Learning Research}, 12, 2825--2830. \\
          \url{https://scikit-learn.org}

    \item Streamlit Inc. (2023). \textit{Streamlit --- The fastest way to build and share data apps}. \\
          \url{https://streamlit.io}

    \item Géron, A. (2022). \textit{Hands-On Machine Learning with Scikit-Learn, Keras, and TensorFlow}
          (3rd ed.). O'Reilly Media.

    \item McKinney, W. (2010). Data Structures for Statistical Computing in Python.
          \textit{Proceedings of the 9th Python in Science Conference}, 445, 51--56.
          (Pandas library.) \\
          \url{https://pandas.pydata.org}

    \item Hunter, J. D. (2007). Matplotlib: A 2D Graphics Environment.
          \textit{Computing in Science \& Engineering}, 9(3), 90--95. \\
          \url{https://matplotlib.org}
\end{enumerate}

% ═══════════════════════════════════════════════════════════════════════════════
\end{document}
